\documentclass[10pt,journal]{IEEEtran}

\usepackage[T1]{fontenc}
\usepackage[utf8]{inputenc}
\usepackage{amsmath}
\usepackage{amssymb}
\usepackage{amsfonts}
\usepackage{mathtools}
\usepackage{graphicx}
\usepackage{booktabs}
\usepackage{array}
\usepackage{float}
\usepackage{caption}
\usepackage{xcolor}
\usepackage{algorithm}
\usepackage{algpseudocode}
\usepackage{cite}
\usepackage{hyperref}
\hypersetup{colorlinks=true, linkcolor=blue, citecolor=blue, urlcolor=blue,
  pdftitle={Vel-v5.5: A Decentralized 3D Volumetric Coverage Protocol under Aerodynamic Wind Perturbations},
  pdfauthor={}}

\newcommand{\vecc}[1]{\vec{\mathbf{#1}}}
\newcommand{\norm}[1]{\left\lVert#1\right\rVert}

\begin{document}

\title{Vel-v5.5: A Decentralized 3D Volumetric Coverage Protocol\\ under Aerodynamic Wind Gust Perturbations}

\author{Manuscript submitted for double-anonymous review to IEEE T-ASE.}

\markboth{Manuscript submitted to IEEE T-ASE, 2026}{Anonymous et al.: Vel-v5.5}

\maketitle

\begin{abstract}
This paper presents Vel-v5.5, a three-dimensional extension of the Vel-v4.3
decentralized swarm coverage protocol, and its empirical characterization
under stochastic aerodynamic wind gust perturbations. Vel-v5.5 retains the
core Map-Aware Gradient Repulsion (MAGR) steering law validated in Vel-v4.3
-- each agent scans a local Moore neighborhood of a shared occupancy grid and
steers directly toward the nearest unexplored cell -- and extends it from a
two-dimensional cell grid to a three-dimensional voxel grid, adding
gravitational drift, stochastic wind gust perturbation, and elastic boundary
collision physics appropriate to aerial multi-rotor deployment. We report a
full factorial sweep of wind gust intensity $G_{\text{wind}} \in \{0, 2, 4,
6, 8\}$ at fixed population $N=150$, momentum $\mu=0.99$, MAGR steering
factor $G=1.5$, and scan radius $r=5$, comprising 2{,}500 independent trials
(500 per gust level) on a $32\times32\times32$ voxel grid. Results show that
zero-gust operation collapses spatio-temporal efficiency to
$\bar{\eta}=0.236$ ($\sigma=0.0059$), that efficiency rises sharply under
mild-to-moderate gust and peaks at $G_{\text{wind}}=4$
($\bar{\eta}=0.666$, $\sigma=0.0034$), and that efficiency declines modestly
under high-intensity gust ($\bar{\eta}=0.637$ at $G_{\text{wind}}=8$). We
interpret the zero-gust collapse as evidence that deterministic frontier
steering alone is prone to trapping agents in locally stable, low-coverage
configurations, and that stochastic wind perturbation functions as an
exploration-enhancing noise source rather than purely as a disturbance to be
rejected. We report this protocol and its empirical characterization as an
engineering specification, and explicitly identify the limits of what is
and is not established by the present study: no formal convergence
guarantee is provided for the discrete steering law, no baseline comparison
against classical potential-field or frontier-exploration methods has yet
been run under identical conditions, and no physical hardware validation
has been performed.
\end{abstract}

\begin{IEEEkeywords}
Swarm robotics, decentralized coordination, occupancy grid, map-aware
repulsion, spatio-temporal efficiency, aerodynamic disturbance rejection,
coverage path planning, multi-rotor systems.
\end{IEEEkeywords}

\noindent\textit{Note to Practitioners---}This work addresses a practical
problem in deploying groups of small flying robots (drones) to search or map
a three-dimensional space, such as a warehouse interior, a collapsed
building, or a stretch of airspace, without relying on GPS or a central
controller telling each drone where to go. Each drone instead follows a
simple local rule: check a shared digital map for the nearest not-yet-visited
spot nearby, and fly toward it. We tested this rule in simulation under
different levels of simulated wind, running 2,500 test flights. The result
may be counterintuitive to practitioners used to treating wind as purely a
problem to be corrected for: completely still air produced the worst
coverage of any condition tested, while moderate wind produced the best
coverage. This suggests that for this type of decentralized, locally-reactive
drone behavior, some ambient disturbance is actually helpful, and that
operators deploying this kind of system in unusually calm conditions (e.g.,
indoors) may want to deliberately add a small amount of randomness to each
drone's flight path rather than flying a perfectly steady, deterministic
route. This paper reports simulation results only; the specific numeric
recommendations have not yet been validated on physical hardware, nor
compared side-by-side against other existing drone-swarm coordination
methods, and practitioners should treat the findings as an initial
characterization rather than a final, field-ready specification.

\section{Introduction}
\label{sec:intro}

Autonomous volumetric coverage using decentralized multi-agent robotic
swarms is a capability of growing importance across planetary surface
mapping, disaster search-and-rescue, atmospheric plume tracking, and
industrial inspection \cite{latombe1991, brambilla2013review}. A recurring
obstacle in decentralized coverage is the \textit{Redundant Search
Problem}: absent shared spatial memory or a global coordinator, independent
agents repeatedly re-traverse already-explored regions, degrading coverage
efficiency as the mission progresses.

The Vel protocol family addresses this problem through a two-layer
decentralized architecture: a Shared Spatial Memory (SSM) layer, in which
every agent reads and writes to a common binary occupancy grid, and a
Kinetic Vector Logic (KVL) layer, in which each agent's velocity is updated
by a local steering law computed from that shared map. Vel-v4.3
\cite{velv43} introduced
Map-Aware Gradient Repulsion (MAGR): rather than steering outward from the
swarm's geometric center (a design that produces a documented
\textit{boundary pile-up} failure mode as agent population scales), each
agent scans a local neighborhood of the SSM and steers directly toward the
nearest unexplored cell. Across 105{,}000 simulated trials, Vel-v4.3
established $\mu=0.99$ and a population-dependent MAGR steering factor
($G=1.5$ at $N=300$, $G=2.75$ at $N=150$) as empirical optima on a
two-dimensional $50\times50$ grid.

Vel-v5.5 extends this validated 2D algorithm into a three-dimensional voxel
environment appropriate for aerial multi-rotor deployment, and asks a
question Vel-v4.3 could not: how does the MAGR frontier-steering law behave
under realistic aerodynamic disturbance? We add three physical elements
absent from Vel-v4.3's abstract 2D grid: a constant gravitational drift
term, a stochastic wind gust perturbation applied to agent velocity at each
timestep, and elastic boundary collision physics in place of a simple
positional clip. We then run a full factorial sweep of wind gust intensity
and report the resulting efficiency characteristics.

This paper supersedes an earlier draft of Vel-v5.0 that formulated MAGR as
a continuous potential field (agent repulsion, boundary distance gradient,
and exploration attraction terms summed as a single steering force) and
paired this formulation with a Lyapunov-based asymptotic convergence proof.
That earlier draft's reference simulation implementation did not, in fact,
apply the derived steering force to agent velocity at all -- an
implementation/theory mismatch discovered during manuscript review. Rather
than retrofit a proof to an untested formulation, Vel-v5.5 instead extends
the simpler, already-validated Vel-v4.3 steering law into 3D and reports
its behavior purely empirically. We regard this as the more defensible
path: an engineering protocol whose claims are limited to what has actually
been simulated and measured.

The remainder of this paper is organized as follows. Section~\ref{sec:related}
reviews related work. Section~\ref{sec:architecture} describes the Vel-v5.5
system architecture. Section~\ref{sec:metrics} defines the performance
metrics used throughout. Section~\ref{sec:study} describes the wind gust
sweep methodology. Section~\ref{sec:results} reports and analyzes the
empirical results. Section~\ref{sec:mvr} provides a strategic
implementation protocol for physical deployment. Section~\ref{sec:limits}
states the explicit limitations of the present study, and
Section~\ref{sec:conclusion} concludes.

\section{Related Work}
\label{sec:related}

\subsection{Coverage Path Planning}
Coverage Path Planning (CPP) is the task of determining a path that passes
over all points of an area or volume of interest while avoiding obstacles
\cite{yamauchi1997}. A considerable body of research has addressed the CPP
problem across domains including vacuum cleaning robots, autonomous
underwater vehicles, demining robots, and aerial survey platforms. Classical
CPP approaches such as boustrophedon decomposition partition the
environment into cells and plan optimal traversal sequences for a single
agent, but do not natively address the coordination challenges that arise
when multiple agents must cover the same volume without collision or
redundant re-traversal.

\subsection{Ant Colony Optimization}
Ant Colony Optimization (ACO) models the pheromone-trail communication of
biological ants to solve combinatorial optimization problems, and has been
extensively applied to multi-robot path planning, where pheromone
concentration guides agents toward unexplored regions \cite{dorigo2006ant}.
The key limitation of ACO for coverage tasks is stagnation: once pheromone
trails converge on a particular traversal pattern, agents are difficult to
redirect toward remaining unexplored regions, and the approach requires
careful tuning of pheromone decay parameters to remain adaptive. Vel
differs from ACO in a fundamental architectural choice: rather than a
continuous pheromone concentration field, Vel's SSM layer uses a binary
occupancy grid in which cells are either explored or unexplored with no
intermediate state. This eliminates the stagnation problem by construction
-- a voxel is explored or it is not -- and removes the need for pheromone
decay tuning entirely.

\subsection{Particle Swarm Optimization}
Particle Swarm Optimization (PSO) models the collective motion of flocking
agents to optimize continuous objective functions, and has been adapted for
multi-robot coverage tasks in which agents are attracted toward unexplored
regions by a shared global-best position \cite{banks2007}. PSO-based
coverage systems are vulnerable to premature convergence: once the global
best position is identified, all agents converge toward it rather than
distributing across the remaining unexplored volume. Vel's KVL layer draws
conceptual inspiration from PSO's velocity update rule -- specifically the
momentum term $\mu$ that damps high-frequency direction changes -- but
replaces PSO's global attractor with a map-aware local gradient, avoiding
premature convergence by ensuring each agent steers toward the nearest
locally unexplored voxel rather than a globally shared target.

\subsection{Centroidal Voronoi Tessellation and Lloyd's Algorithm}
Centroidal Voronoi Tessellations (CVT), computed via Lloyd's algorithm
\cite{lloyd1982}, partition a coverage domain into cells such that each
agent is responsible for the region closest to it, and have been formally
extended to mobile sensing network coverage control \cite{cortes2004}.
Extending CVT-based coverage from 2D to 3D volumes introduces substantial
computational and geometric complexity: computing 3D Voronoi polyhedra in
real time requires continuous, high-precision positioning data from every
neighboring agent, demanding communication bandwidth and onboard processing
power that are often unavailable on resource-constrained micro-aerial
vehicle (MAV) platforms. Furthermore, classical CVT formulations are
static and do not naturally adapt to an actively evolving spatial priority
field as regions are progressively mapped. Vel-v5.5 achieves emergent
volumetric coverage through localized, coordinate-blind gradient descent
on a binary voxel grid, avoiding both the communication overhead and the
static-partition limitation of CVT-based approaches.

\subsection{Frontier-Based Exploration}
Frontier-based exploration architectures identify the boundary between
mapped free space and unmapped space, directing agents toward these
``frontiers'' to expand the map \cite{yamauchi1997}. Classical frontier
detection in 3D is typically a centralized process requiring global map
synchronization and path-planning algorithms such as $A^*$ search or Fast
Marching Methods to compute collision-free trajectories to identified
frontiers -- a bottleneck that scales poorly with agent count and is
vulnerable to single-point-of-failure network drops. Vel's MAGR steering
law achieves the target-attraction behavior of frontier-based schemes
through a purely localized, reactive scan-and-steer computation, avoiding
the need for global trajectory planning or explicit frontier tracking.

\subsection{Decentralized Multi-Robot Coverage}
Decentralized multi-robot coverage architectures achieve competitive
coverage efficiency while eliminating the single point of failure inherent
to centralized command structures \cite{brambilla2013review}.
Potential-field-based decentralized swarms, however, are documented to
suffer severe spatial instabilities when navigating near physical
boundaries or non-convex obstacles \cite{saska2014cooperative}. Vel-v4.3's
introduction of MAGR directly addressed an analogous instability (boundary
pile-up under center-outward repulsion) in a 2D setting by modeling the
physical perimeter through local map awareness rather than global geometric
bearing. The present study evaluates whether this same architectural
solution remains effective when the steering law is extended into three
dimensions and subjected to physical aerodynamic disturbance.

\subsection{Occupancy Grid Representations}
The occupancy grid representation underlying Vel's SSM layer originates in
probabilistic mobile robot perception and mapping research, where
individual cells encode the probability of physical occupancy under sensor
uncertainty. Vel simplifies this representation to a deterministic binary
state (explored/unexplored), appropriate for a coverage task in which the
quantity of interest is certainty of prior visitation rather than
probability of physical occupancy.

\subsection{Positioning of Vel-v5.5}
Vel-v5.5 occupies a distinct position in the swarm coverage landscape.
Unlike ACO, it uses deterministic binary memory rather than probabilistic
pheromone concentration. Unlike PSO, it uses local map-aware gradient
descent rather than global positional attraction. Unlike CVT-based
methods, it requires no continuous inter-agent positioning communication
and adapts naturally as the occupancy field evolves. Unlike classical
frontier-exploration CPP methods, it is fully decentralized with no
pre-planned or centrally computed traversal sequence. The present
contribution is the extension of this already-validated 2D steering law
into a three-dimensional voxel environment and its systematic
characterization under stochastic aerodynamic disturbance -- a condition
not modeled in any of the coverage architectures reviewed above.

\section{System Architecture}
\label{sec:architecture}

Vel-v5.5 retains Vel-v4.3's two-layer architecture.

\begin{figure}[t]
\centering
\fbox{\parbox{0.92\columnwidth}{\centering
\textbf{Vel-v5.5 Two-Layer Architecture} \\[0.3cm]
\fbox{\parbox{0.85\columnwidth}{\centering \textbf{Shared Spatial Memory (SSM)} \\ $32\times32\times32$ binary voxel occupancy grid \\ read/write by all agents every timestep}} \\[0.2cm]
$\updownarrow$ \\[0.2cm]
\fbox{\parbox{0.85\columnwidth}{\centering \textbf{Kinetic Vector Logic (KVL)} \\ Moore-neighborhood scan $\rightarrow$ MAGR steering \\ $+$ momentum $+$ gravity drift $+$ wind gust \\ $+$ elastic boundary collision}} \\[0.2cm]
$\downarrow$ \\[0.2cm]
\fbox{\parbox{0.6\columnwidth}{\centering Agent position \& velocity update}}
}}
\caption{Two-layer Vel-v5.5 architecture. Every agent reads the shared SSM
to locate its steering target, then updates its own velocity through the
KVL layer independently -- no direct agent-to-agent communication channel
is required.}
\label{fig:architecture}
\end{figure}

Fig.~\ref{fig:architecture} summarizes the flow of information between the
two layers, described in detail below.

\subsection{Shared Spatial Memory (SSM)}
The operational volume is discretized into a voxel grid $\mathcal{V}$ of
dimensions $32\times32\times32$ (32{,}768 voxels). Each voxel
$c(x,y,z)\in\mathcal{V}$ is a binary state variable:
\begin{equation}
c(x,y,z) = \begin{cases} 0 & \text{Unexplored} \\ 1 & \text{Explored} \end{cases}
\end{equation}
Every agent reads and writes the same SSM, providing implicit coordination
without direct agent-to-agent communication.

\subsection{Kinetic Vector Logic (KVL) and MAGR Steering}
At each timestep, agent $i$ at position $\mathbf{p}_i(t)\in\mathbb{R}^3$
first marks its current voxel explored, then scans the Moore neighborhood
of radius $r$ around that voxel for the nearest unexplored voxel. If one is
found at relative offset $(dx^*,dy^*,dz^*)$, the agent's steering
acceleration is:
\begin{equation}
\vec{R}_G = G \cdot \frac{(dx^*,dy^*,dz^*)}{\norm{(dx^*,dy^*,dz^*)}} \cdot k
\label{eq:magr_found}
\end{equation}
where $G\in\mathbb{R}^+$ is the tunable steering factor and $k=0.12$ is the
fixed scale constant carried over from Vel-v4.3. If no unexplored voxel is
found within radius $r$, the agent falls back to steering directly away
from the swarm's instantaneous geometric center $\bar{\mathbf{p}}(t)$:
\begin{equation}
\vec{R}_G = G \cdot \frac{\mathbf{p}_i(t) - \bar{\mathbf{p}}(t)}{\norm{\mathbf{p}_i(t) - \bar{\mathbf{p}}(t)}} \cdot k
\label{eq:magr_fallback}
\end{equation}

The full velocity update integrates this steering term with momentum,
gravitational drift, and stochastic wind gust perturbation:
\begin{equation}
\mathbf{v}_i(t+1) = \mu\,\mathbf{v}_i(t) + \vec{R}_G - \mathbf{g}_{\text{drift}} + \mathbf{w}_{\text{gust}}(t)
\label{eq:vel_update}
\end{equation}
\begin{equation}
\mathbf{p}_i(t+1) = \mathbf{p}_i(t) + \mathbf{v}_i(t+1)
\end{equation}
where $\mu\in[0,1)$ is the momentum coefficient, $\mathbf{g}_{\text{drift}} =
[0, g_{\text{aero}}, 0]^\top$ is a constant downward gravitational drift
($g_{\text{aero}}=0.082$), and $\mathbf{w}_{\text{gust}}(t)$ is a stochastic
wind gust vector with independent components:
\begin{equation}
w_{j,i}(t) \sim \mathcal{U}(-0.5, 0.5) \cdot G_{\text{wind}}, \quad \forall j\in\{x,y,z\}
\label{eq:gust}
\end{equation}
parameterized by gust intensity $G_{\text{wind}}$. Note that
$G_{\text{wind}}$ (wind gust intensity) and $G$ (MAGR steering factor, Eq.
\ref{eq:magr_found}) are distinct quantities and are varied independently
in this study; $G$ is held fixed at the Vel-v4.3 empirical optimum
throughout, while $G_{\text{wind}}$ is the swept experimental variable.

\subsection{Computational Complexity of 3D Neighborhood Scanning}
Extending MAGR's local neighborhood scan from a 2D Moore neighborhood to a
3D Moore neighborhood changes its computational cost scaling. In Vel-v4.3's
2D formulation, scanning a neighborhood of radius $r$ requires
$(2r+1)^2$ SSM reads per agent per timestep. In the 3D voxel formulation
used here, the equivalent scan requires $(2r+1)^3$ reads per agent per
timestep -- a cubic rather than quadratic growth in $r$. At the scan
radius used throughout this study ($r=5$), this corresponds to 121 reads
per agent per step in 2D versus 1{,}331 reads per agent per step in 3D, an
order-of-magnitude increase in onboard scan cost for the same radius. This
has a direct engineering implication: Vel-v4.3's Study V finding that MAGR
efficiency is radius-invariant in 2D (permitting a small, computationally
cheap scan radius) cannot be assumed to transfer to 3D without
re-verification, since the cubic cost scaling makes the choice of scan
radius a substantially more consequential deployment decision in a voxel
environment than in a 2D cell grid. Re-sweeping scan radius in 3D to locate
the analogous radius-invariance threshold is identified as future work in
Section~\ref{sec:limits}.

\subsection{Elastic Boundary Collision}
When an agent's predicted position exceeds the grid boundary $[0,31]^3$,
its position is capped at the boundary and its velocity along the violated
axis is reflected with restitution coefficient $\epsilon=-0.45$:
\begin{equation}
v_{j,i}(t+1) \leftarrow \epsilon \cdot v_{j,i}(t+1)
\end{equation}

\begin{algorithm}[t]
\caption{Vel-v5.5 Per-Agent Update}
\label{alg:vel55}
\begin{algorithmic}[1]
\Procedure{UpdateAgent}{$i$}
    \State $(gx,gy,gz) \gets \lfloor \mathbf{p}_i \rfloor$
    \State $c(gx,gy,gz) \gets 1$
    \State Scan Moore neighborhood of radius $r$ for nearest unexplored voxel
    \If{found at $(dx^*,dy^*,dz^*)$}
        \State $\vec{R}_G \gets$ Eq.~\ref{eq:magr_found}
    \Else
        \State $\vec{R}_G \gets$ Eq.~\ref{eq:magr_fallback}
    \EndIf
    \State $\mathbf{w}_{\text{gust}} \gets$ sample per Eq.~\ref{eq:gust}
    \State $\mathbf{v}_i \gets \mu\mathbf{v}_i + \vec{R}_G - \mathbf{g}_{\text{drift}} + \mathbf{w}_{\text{gust}}$
    \State $\mathbf{p}_i \gets \mathbf{p}_i + \mathbf{v}_i$
    \State Apply elastic boundary collision if $\mathbf{p}_i$ exceeds $[0,31]^3$
    \State Mark voxel at $\lfloor \mathbf{p}_i \rfloor$ explored if unexplored
\EndProcedure
\end{algorithmic}
\end{algorithm}

\section{Performance Metrics}
\label{sec:metrics}

\subsection{Spatio-Temporal Efficiency}
Consistent with Vel-v4.3, the primary metric is:
\begin{equation}
\eta = \frac{\Delta C}{\Delta T \cdot N}
\end{equation}
where $\Delta C$ is the number of newly explored voxels over the trial,
$\Delta T$ is the number of elapsed timesteps, and $N$ is agent count. This
is equivalently computed here as $\text{Covered}/(\text{Steps}\cdot N)$ over
a fixed 150-step trial.

\subsection{Stability}
Trial-to-trial standard deviation of $\eta$ within a configuration:
\begin{equation}
\sigma = \sqrt{\frac{1}{K}\sum_{k=1}^{K}(\eta_k - \bar{\eta})^2}
\end{equation}

\section{Wind Gust Sweep: Methodology}
\label{sec:study}

We fix population $N=150$, momentum $\mu=0.99$ (the Vel-v4.3 Study IV
optimum), MAGR steering factor $G=1.5$ (the Vel-v4.3 Study III optimum at
$N=300$; re-validating this choice at $N=150$ in 3D is identified as future
work in Section~\ref{sec:limits}), and scan radius $r=5$. We sweep wind
gust intensity $G_{\text{wind}} \in \{0,2,4,6,8\}$, running 500 independent
trials per gust level (2{,}500 trials total) on the $32\times32\times32$
voxel grid described in Section~\ref{sec:architecture}. Each trial runs for
a fixed 150 timesteps. All trials use independent random seeds drawn from a
single seed sequence for reproducibility.

\section{Results}
\label{sec:results}

\begin{table}[t]
\centering
\caption{Wind Gust Sweep Results ($N=150$, $\mu=0.99$, $G=1.5$, $r=5$, 500 trials per level)}
\label{tab:gust_sweep}
\begin{tabular}{@{}ccccc@{}}
\toprule
$G_{\text{wind}}$ & $\bar{\eta}$ & $\sigma_\eta$ & min & max \\
\midrule
0 & 0.2361 & 0.0059 & 0.2197 & 0.2523 \\
2 & 0.6204 & 0.0053 & 0.6036 & 0.6347 \\
4 & \textbf{0.6664} & 0.0034 & 0.6566 & 0.6763 \\
6 & 0.6565 & 0.0029 & 0.6478 & 0.6658 \\
8 & 0.6373 & 0.0028 & 0.6291 & 0.6454 \\
\bottomrule
\end{tabular}
\end{table}

\begin{figure}[t]
\centering
\includegraphics[width=0.95\columnwidth]{figures/fig1_efficiency_vs_gust.png}
\caption{Mean spatio-temporal efficiency vs. wind gust intensity, with
$\pm 1\sigma$ error bars ($n=500$ per level). Efficiency peaks at
$G_{\text{wind}}=4$.}
\label{fig:eff_vs_gust}
\end{figure}

\begin{figure}[t]
\centering
\includegraphics[width=0.95\columnwidth]{figures/fig2_efficiency_boxplot.png}
\caption{Full efficiency distribution by gust level ($n=500$ per level).
Distributions are tight and non-overlapping between the zero-gust and
all wind-perturbed conditions.}
\label{fig:eff_boxplot}
\end{figure}

\subsection{Zero-Gust Collapse}
The zero-gust condition ($G_{\text{wind}}=0$) produces the lowest observed
efficiency, $\bar{\eta}=0.236$, less than half the efficiency of every
wind-perturbed condition tested. Given the tight standard deviation at this
level ($\sigma=0.0059$) and the complete non-overlap of its distribution
with any other tested gust level (Fig.~\ref{fig:eff_boxplot}), this is not
a marginal effect. We attribute this to the deterministic nature of the
nearest-cell steering law under zero disturbance: absent any stochastic
perturbation, agents following identical local scan-and-steer logic from
similar starting conditions can settle into mutually reinforcing,
locally-stable trajectory patterns that repeatedly re-scan the same
already-explored local neighborhoods rather than escaping toward more
distant unexplored volume. This is consistent with known local-minima
trapping behavior in deterministic potential-field and frontier-following
navigation methods generally.

\subsection{Non-Monotonic Response to Gust Intensity}
Efficiency rises sharply from $G_{\text{wind}}=0$ to $G_{\text{wind}}=2$,
continues to a peak at $G_{\text{wind}}=4$ ($\bar{\eta}=0.6664$), then
declines modestly through $G_{\text{wind}}=6$ and $G_{\text{wind}}=8$
(Fig.~\ref{fig:eff_vs_gust}). This shape is consistent with wind gust
functioning as an exploration-enhancing stochastic perturbation at moderate
intensity -- breaking agents out of the local trapping behavior identified
above -- while at higher intensity the perturbation increasingly dominates
and degrades the steering signal itself, consistent with the elastic
boundary collision mechanism dissipating a larger share of trajectory
energy as gust-driven boundary impacts become more frequent.

\subsection{Stability Decreases with Distance from Zero-Gust}
Standard deviation is highest at $G_{\text{wind}}=0$ ($\sigma=0.0059$) and
generally decreases as gust intensity rises to $G_{\text{wind}}=6$
($\sigma=0.0029$) and $G_{\text{wind}}=8$ ($\sigma=0.0028$), with
$G_{\text{wind}}=2$ as a mild exception ($\sigma=0.0053$). This suggests
that moderate-to-high wind gust, in addition to raising mean efficiency
relative to the zero-gust case, also produces more trial-to-trial
consistent behavior -- a potentially useful property for deployment
scenarios requiring predictable coverage completion times.

\subsection{Statistical Significance}
\label{sec:stats}
To confirm that the differences reported above reflect a genuine effect
rather than sampling noise, we conducted pairwise Welch's $t$-tests between
each adjacent gust level in Table~\ref{tab:gust_sweep}, with Cohen's $d$
computed as the standardized effect size. Results are summarized in
Table~\ref{tab:ttests}.

\begin{table}[t]
\centering
\caption{Pairwise Welch's $t$-tests Between Adjacent Gust Levels ($N=150$ dataset)}
\label{tab:ttests}
\begin{tabular}{@{}ccccc@{}}
\toprule
Comparison & $t$ & $p$ & Cohen's $d$ \\
\midrule
$G_{\text{wind}}{=}0$ vs.\ $2$ & $-1080.5$ & $<10^{-300}$ & $68.4$ \\
$G_{\text{wind}}{=}2$ vs.\ $4$ & $-163.4$ & $<10^{-300}$ & $10.3$ \\
$G_{\text{wind}}{=}4$ vs.\ $6$ & $49.6$ & $\approx10^{-268}$ & $-3.1$ \\
$G_{\text{wind}}{=}6$ vs.\ $8$ & $107.4$ & $<10^{-300}$ & $-6.8$ \\
\bottomrule
\end{tabular}
\end{table}

Every adjacent-level comparison is statistically significant at an
extreme margin, and every effect size exceeds conventional thresholds for
a ``large'' effect ($|d|>0.8$) by more than an order of magnitude. The
zero-gust-to-moderate-gust transition in particular ($d=68.4$) is among
the largest effect sizes we are aware of being reported in a swarm
coverage parameter study, reflecting the near-total non-overlap of the
zero-gust distribution with every wind-perturbed condition visible in
Fig.~\ref{fig:eff_boxplot}. We conclude that the non-monotonic response to
gust intensity reported in this paper is a robust feature of the tested
configuration and not an artifact of sampling variation.

\subsection{Population Density Effect: Preliminary N=300 Comparison}
Vel-v4.3's own Study I identified a genuine diminishing-returns effect in
Coordination Quality ($\eta/N$) as population scales, driven by
scan-overlap congestion. To check whether an analogous effect appears in
the 3D setting, we ran a preliminary comparison at $N=300$ (double the
main study's population) at three representative gust levels
($G_{\text{wind}} \in \{0,4,8\}$), with 20 trials per level -- a smaller
sample than the main study's 500 trials per level, run as a pilot check
rather than a full characterization. Results are shown in
Table~\ref{tab:n300} and Fig.~\ref{fig:n150_vs_n300}.

\begin{table}[t]
\centering
\caption{N=300 Preliminary Comparison (20 trials per gust level)}
\label{tab:n300}
\begin{tabular}{@{}cccc@{}}
\toprule
$G_{\text{wind}}$ & $\bar{\eta}$ ($N=300$) & $\bar{\eta}$ ($N=150$) & Ratio \\
\midrule
0 & 0.1804 & 0.2361 & 0.76 \\
4 & 0.5014 & 0.6664 & 0.75 \\
8 & 0.4776 & 0.6373 & 0.75 \\
\bottomrule
\end{tabular}
\end{table}

\begin{figure}[t]
\centering
\includegraphics[width=0.95\columnwidth]{figures/fig3_n150_vs_n300.png}
\caption{Efficiency vs. gust intensity at $N=150$ (full study, $n=500$ per
level) and $N=300$ (preliminary comparison, $n=20$ per level). The
qualitative pattern -- zero-gust collapse, peak at moderate gust -- holds
at both population sizes, though absolute efficiency is systematically
lower at $N=300$.}
\label{fig:n150_vs_n300}
\end{figure}

At all three tested gust levels, doubling population from 150 to 300
reduces mean efficiency by a remarkably consistent ratio of approximately
0.75--0.76, while the qualitative shape of the gust-response curve --
zero-gust collapse, peak at moderate gust, mild decline at high gust --
is preserved at both population sizes. This is consistent with the same
scan-overlap congestion mechanism Vel-v4.3 documented in 2D: because
efficiency $\eta$ is normalized by $N$, a larger swarm operating in the
same volume must achieve proportionally more coverage to maintain
constant $\eta$, and increased agent density raises the likelihood that
multiple agents simultaneously target the same nearby unexplored voxel.
We emphasize that this comparison uses a substantially smaller sample
than the main $N=150$ study and should be treated as a preliminary
finding pending a full $N=300$ sweep at 500 trials per level, consistent
with Vel-v4.3's own systematic population-sweep methodology.

\section{Strategic Implementation Protocol (MVR)}
\label{sec:mvr}

Based on the results in Section~\ref{sec:results} and the Vel-v4.3
findings this protocol extends, we identify the following Minimum Viable
Requirements for physical multi-rotor deployment:

\begin{enumerate}
\item \textbf{MVR-1: Controller Update Frequency $\geq 50$ Hz.} Consistent
with Vel-v4.3's KVL update rate requirement.
\item \textbf{MVR-2: Mesh Synchronization Latency $< 15$ ms.} To limit
Ghost Redundancy (multiple agents targeting the same unexplored voxel
before SSM writes propagate).
\item \textbf{MVR-3: Momentum Coefficient $\mu \geq 0.95$.} Carried over
from Vel-v4.3 Study IV; not independently re-swept in the present 3D study
(see Section~\ref{sec:limits}).
\item \textbf{MVR-4: Do not operate in still-air conditions.} Section
\ref{sec:results} shows zero-gust operation produces the lowest observed
efficiency of any tested condition. Where natural wind is unavailable or
insufficient, an artificial stochastic perturbation to agent velocity is
recommended to replicate the exploration-enhancing effect observed here.
\item \textbf{MVR-5: Target Operating Gust Range $G_{\text{wind}} \in
[3,5]$.} Bracketing the empirical peak at $G_{\text{wind}}=4$.
\end{enumerate}

\section{Discussion}
\label{sec:discussion}

\subsection{Noise-Assisted Exploration: A Broader Context}
The finding that stochastic perturbation improves a deterministic search
process is not unique to swarm coverage. In combinatorial optimization,
simulated annealing deliberately introduces a temperature-controlled
stochastic perturbation to escape local minima that a purely greedy search
would remain trapped in \cite{kirkpatrick1983}. In nonlinear dynamical
systems, stochastic resonance describes conditions under which added noise
improves detection or response of an otherwise sub-threshold signal. The
zero-gust collapse and subsequent moderate-gust peak reported in
Section~\ref{sec:results} are broadly consistent with this class of
phenomena: the MAGR steering law, as a purely deterministic
nearest-cell-seeking rule, appears to define a search landscape containing
locally stable configurations from which a swarm of identically-programmed
agents cannot escape without external perturbation. Wind gust, in this
framing, is not merely an environmental disturbance to be rejected by
robust control, but is itself doing useful computational work: it is the
source of the randomness the steering law itself does not provide. This
reframing has a direct engineering implication distinct from the MVR
recommendations in Section~\ref{sec:mvr}: if a deployment environment is
unusually still, deliberately injecting a small stochastic perturbation
into the KVL velocity update (rather than merely tolerating ambient wind)
may be preferable to relying on incidental turbulence.

\subsection{Why Efficiency Declines at High Gust Intensity}
The decline in mean efficiency beyond $G_{\text{wind}}=4$ (Section
\ref{sec:results}) is consistent with two compounding mechanisms. First, as
gust magnitude grows relative to the fixed MAGR steering magnitude ($G
\cdot k$ in Eq.~\ref{eq:magr_found}), the wind term increasingly dominates
the velocity update in Eq.~\ref{eq:vel_update}, degrading the directional
signal that carries an agent toward its identified target voxel. Second,
higher gust intensity increases the frequency with which agents' predicted
positions exceed the grid boundary, triggering the elastic collision
mechanism (Eq.~\ref{eq:vel_update}) more often; because the restitution
coefficient $\epsilon=-0.45$ dissipates a substantial fraction of kinetic
energy on each boundary impact, more frequent boundary collisions convert
a larger share of the agent's total trajectory energy into dissipated
energy rather than productive exploration motion. Distinguishing the
relative contribution of these two mechanisms would require an ablation
in which boundary collision frequency is logged directly -- a
straightforward instrumentation addition identified as future work.

\subsection{Relationship to the Superseded Vel-v5.0 Formulation}
It is worth stating plainly why this paper does not attempt to reconcile
its findings with the continuous potential-field formulation and Lyapunov
convergence proof of the superseded Vel-v5.0 draft (Section
\ref{sec:intro}). The two formulations are different control laws: the
potential-field formulation computes steering as the negative gradient of
a smooth composite potential summing agent-repulsion, boundary-distance,
and exploration-attraction terms, while the discrete formulation studied
here computes steering as a unit vector toward the single nearest
unexplored voxel located by discrete neighborhood search. A convergence
proof derived for the former does not apply to the latter, and we regard
an unsupported claim of formal convergence as a more serious defect in an
engineering protocol paper than the simple absence of one. The empirical
characterization in this paper is offered as a substitute standard of
evidence appropriate to the discrete steering law actually implemented and
tested.

\section{Threats to Validity}
\label{sec:limits}

We organize the limitations of this study using the standard internal,
external, and construct validity framing common to empirical robotics and
software engineering research.

\subsection{Internal Validity}
\textit{Does the reported effect actually stem from gust intensity, rather
than a confound?} All configuration parameters other than
$G_{\text{wind}}$ (population, momentum, MAGR steering factor, scan radius,
grid size, trial length, and random seed sequence) were held fixed across
the sweep, and every trial's random seed is drawn independently from a
common seed sequence (Section~\ref{sec:repro}), so no shared randomness
could produce a spurious cross-condition correlation. The statistical
tests in Section~\ref{sec:stats} confirm the observed differences are
far larger than sampling variation could plausibly produce. We consider
internal validity for the specific claim ``efficiency varies
non-monotonically with $G_{\text{wind}}$ under this configuration'' to be
well supported.

\subsection{External Validity}
\textit{Do these findings generalize beyond the tested configuration?}
This is the weakest aspect of the present study. All results use a single
population size for the main study ($N=150$), a single MAGR steering
factor and momentum coefficient (both carried over from Vel-v4.3's 2D
optimum rather than independently re-derived in 3D), a single scan radius,
and a single grid size and trial length. The preliminary $N=300$ comparison
(Section~\ref{sec:results}) suggests the qualitative gust-response pattern
is not an artifact of the specific $N=150$ configuration, but this was
tested at reduced sample size and only three gust levels. We explicitly do
not claim that the reported optimum ($G_{\text{wind}}=4$) generalizes to
other population sizes, steering factors, momentum coefficients, scan
radii, or grid dimensions without further testing.

\subsection{Construct Validity}
\textit{Does spatio-temporal efficiency $\eta$ actually measure what a
deployment engineer cares about?} $\eta$ is normalized by both elapsed time
and agent count, and is therefore a reasonable proxy for the economic cost
of a coverage mission (fewer agent-timesteps per unit coverage is cheaper
to operate), but it does not directly capture other properties a physical
deployment may require, such as worst-case time-to-completion, energy
consumption per agent, or resilience to individual agent failure. No claim
in this paper should be read as a claim about those other properties.

\subsection{Summary of Untested Claims}
Consolidating the above, we state explicitly what remains untested:

\begin{itemize}
\item \textbf{No formal convergence guarantee.} Unlike a continuous
potential-field formulation, no Lyapunov-style or other formal proof of
asymptotic coverage is provided for the discrete nearest-cell steering law
studied here. The empirical results in Section~\ref{sec:results} should be
read as evidence of practical behavior under the tested conditions, not as
a proven guarantee.
\item \textbf{No baseline comparison.} This study characterizes Vel-v5.5's
own response to wind gust intensity; it does not compare Vel-v5.5 against
classical artificial potential field or frontier-exploration methods run
under identical 3D, wind-perturbed conditions. Such a baseline comparison
is necessary to substantiate any claim of relative advantage over existing
methods and is identified as the immediate next step for this line of work.
\item \textbf{Population size only preliminarily explored.} The main study
uses $N=150$. Section~\ref{sec:results} reports a preliminary $N=300$
comparison (20 trials per level, versus 500 for the main study) that finds
a consistent efficiency reduction ratio and a preserved qualitative
gust-response shape, suggestive of the same scan-overlap congestion
Vel-v4.3 documented in 2D. This finding should be treated as preliminary;
a full $N=300$ sweep at matching sample size (500 trials per level), and
ideally additional population points as in Vel-v4.3's own three-point
sweep, is needed before a population-dependent claim can be made with the
same confidence as the gust-intensity findings.
\item \textbf{MAGR steering factor, momentum, and scan radius not
re-swept in 3D.} $G=1.5$, $\mu=0.99$, and $r=5$ are carried over from
Vel-v4.3's 2D optimum rather than independently re-derived for the 3D
voxel environment. Section~\ref{sec:architecture}'s complexity analysis
indicates scan radius in particular may not transfer directly given the
cubic rather than quadratic cost scaling in 3D.
\item \textbf{Simulation only.} No physical hardware validation has been
performed. The MVR specifications in Section~\ref{sec:mvr} are design
targets, not validated hardware requirements.
\end{itemize}

\section{Reproducibility and Data Availability}
\label{sec:repro}

All 2{,}500 trial records underlying Table~\ref{tab:gust_sweep} and the
20-trial-per-level $N=300$ comparison underlying Table~\ref{tab:n300} were
generated by a single-file Python simulation implementing
Algorithm~\ref{alg:vel55} exactly as specified. Each trial uses an
independently-drawn random seed from a common seed sequence, and raw
per-trial efficiency values (not only summary statistics) are retained in
the accompanying dataset, permitting independent recomputation of every
statistic and figure reported in this paper, including the pairwise
significance tests in Section~\ref{sec:stats}.

\section{Conclusion}
\label{sec:conclusion}

Vel-v5.5 extends the empirically-validated Vel-v4.3 MAGR frontier-steering
protocol from a 2D cell grid into a 3D voxel environment subject to
aerodynamic disturbance. Across 2{,}500 trials sweeping wind gust
intensity, we find that stochastic wind perturbation is not merely a
disturbance to be rejected but functions as an exploration-enhancing
mechanism: zero-gust operation collapses efficiency to less than half of
any wind-perturbed condition tested, efficiency peaks at moderate gust
intensity ($G_{\text{wind}}=4$), and efficiency declines modestly under
higher-intensity gust. These findings are reported strictly as empirical
observations under the tested conditions; Section~\ref{sec:limits}
identifies baseline comparison, population re-sweeping, and formal
theoretical characterization of the discrete steering law as necessary next
steps before broader claims of protocol superiority can be substantiated.

\section*{Acknowledgements}
Acknowledgements are omitted for double-anonymous review and will be
included in the camera-ready version upon acceptance. Generative AI tools
were used for simulation code implementation assistance, LaTeX formatting,
and prose editing; all research concepts, experimental design, and
scientific conclusions are the original work of the author(s).

\begin{thebibliography}{9}

\bibitem{velv43}
Anonymous, ``Vel-v4.3: A Decentralized Swarm Protocol for Autonomous Volume Dominance,'' anonymized for double-blind review, 2026.

\bibitem{latombe1991}
J.-C.~Latombe, \emph{Robot Motion Planning}. Boston, MA: Kluwer Academic Publishers, 1991.

\bibitem{brambilla2013review}
M.~Brambilla, E.~Ferrante, M.~Birattari, and M.~Dorigo, ``Swarm robotics: a review from the swarm engineering perspective,'' \emph{Swarm Intelligence}, vol.~7, no.~1, pp. 1--41, 2013.

\bibitem{yamauchi1997}
B.~Yamauchi, ``A frontier-based approach for autonomous exploration,'' in \emph{Proc. IEEE Int. Symp. Comput. Intell. Robot. Autom. (CIRA)}, 1997, pp. 146--151.

\bibitem{saska2014cooperative}
M.~Saska et al., ``Autonomous deployment of swarms of micro-aerial vehicles in cooperative surveillance,'' in \emph{Proc. ICUAS}, 2014.

\bibitem{dorigo2006ant}
M.~Dorigo, M.~Birattari, and T.~Stutzle, ``Ant colony optimization,'' \emph{IEEE Comput. Intell. Mag.}, vol.~1, no.~4, pp. 28--39, 2006.

\bibitem{banks2007}
A.~Banks, J.~Vincent, and C.~Anyakoha, ``A review of particle swarm optimization. Part I: background and development,'' \emph{Nat. Comput.}, vol.~6, no.~4, pp. 467--484, 2007.

\bibitem{lloyd1982}
S.~Lloyd, ``Least squares quantization in {PCM},'' \emph{IEEE Trans. Inf. Theory}, vol.~28, no.~2, pp. 129--137, 1982.

\bibitem{cortes2004}
J.~Cortés, S.~Martínez, T.~Karatas, and F.~Bullo, ``Coverage control for mobile sensing networks,'' \emph{IEEE Trans. Robot. Autom.}, vol.~20, no.~2, pp. 243--255, 2004.

\bibitem{kirkpatrick1983}
S.~Kirkpatrick, C.~D.~Gelatt, and M.~P.~Vecchi, ``Optimization by simulated annealing,'' \emph{Science}, vol.~220, no.~4598, pp. 671--680, 1983.

\end{thebibliography}

\end{document}
